\chapter[Effectivizing GMT]{Applications to Geometric Measure Theory}\label{chapter:effectivizing_GMT}

\section{Two Approaches to the Combinatorics of Geometric Measure Theory} \label{section:Effective_versions_of_Combinatorial_Proofs}

We continue the discussion from Section~\ref{section:Algorithmic_Randomness} on $(C, \delta, s)$-sets and their relation with Hausdorff dimension. To set up our discussion, let us recall more notation from \cite{Mat95}. For any natural numbers $0 < n < m$, the notation $G(m, n)$ refers to the \textit{Grassmannian manifold} of all $n$-dimensional subspaces of $\mathbb{R}^m$. Whenever $V \in G(m, n)$, let $\pi_V : \mathbb{R}^m \to V$ denote to the orthogonal projection onto $V$. We take the usual operator norm $\norm{\cdot}$ for linear maps as our metric on $G(m, n)$. Let $\mathcal{O}(m)$ denote the \textit{$m$-dimensional orthogonal group}, which consists of all linear maps $g : \mathbb{R}^m \to \mathbb{R}^m$ preserving the inner product in $\mathbb{R}^m$. Considering the \textit{Haar measure} $\theta_m$ on $\mathcal{O}(m)$, which is an invariant probability measure uniformly distributed on $\mathcal{O}(m)$, then $\gamma_{m, n}$ will denote the unique orthogonally-invariant, Radon probability measure on $G(m, n)$ arising from this Haar measure. That is, for any $X \subseteq G(m, n)$ and any fixed $V_0 \in X$,
\begin{align*}
    \gamma_{m, n}(X) :=  \theta_m(\set{g \in \mathcal{O}(m)}{g(V_0) \in X}).
\end{align*}
See Section 3.9 in \cite{Mat95} for more details. One celebrated fact about how Hausdorff dimension behaves under orthogonal projections onto subspaces of $\mathbb{R}^m$ is commonly known as \textit{the Marstrand-Mattila Projection Theorem}.

% {\color{red} Discuss briefly the history of Marstrand-Mattila (planar to general Euclidean space), then its widening to sets whose Hausdorff and packing dimensions agree, etc.}
\begin{theorem}[Marstrand-Mattila Projection Theorem; c.f. Theorem 5.8 in \cite{Mat15}]
    \label{theorem:marstrand_projection_original}
    Let $0 < n < m$ and $X \subseteq \mathbb{R}^m$ be analytic. Then,
    \begin{align*}
        \dim_{\operatorname{H}} \pi_V(X) = \min \seq{\dim_{\operatorname{H}} X, n} \qquad \text{for $\gamma_{m, n}$-almost all $V \in G(m, n)$}.
    \end{align*}
\end{theorem}

The $n = 1$ case of Theorem~\ref{theorem:marstrand_projection_original}--as well as other basic facts about how Hausdorff dimension behaves under translations, intersections, and Cartesian products--now have alternative, effective proofs involving an application of the Point-to-Set Principle \cite{LS17, Lut21}. Such effective proofs offer the advantage of refining what ``almost all'' means usually by finding a sufficiently strong notion of algorithmic randomness which witnesses the theorem.

In particular, N.~Lutz and D.~Stull in \cite{LS17} found an effective proof that the conclusion of the $n = 1$ case of the Marstrand-Mattila Projection Theorem still holds under the alternate hypothesis that $\dim_{\operatorname{H}} X = \dim_{\operatorname{P}} X$. Shortly after, D. Stull extended this conclusion for a broader class of sets which are said to possess \textit{optimal oracles} \cite{Stu21}, a condition which has not yet been fully characterized in non-computability theoretic terms. 

However, there was a non-computability theoretic response to the version of the Marstrand-Mattila Projection Theorem under the condition that Hausdorff and packing dimensions agree. Namely, T.~Orponen in \cite{Orp21} employed combinatorial methods to extend the result to all $0 < n < m$ (see Theorem 1.2 in \cite{Orp21}). 

He further proved that a weaker conclusion holds for \textit{any} subset of $\mathbb{R}^m$.
\begin{theorem}[Theorem 1.3 in \cite{Orp21}]
    \label{theorem:Orponen_weak_Marstrand}
    Let $0 < n < m$ and $X \subseteq \mathbb{R}^m$. Then 
    \begin{align*}
        \dim_{\operatorname{P}} \pi_V(X) = \min \seq{\dim_{\operatorname{H}} X, n} \qquad \text{for $\gamma_{m, n}$-almost all $V \in G(m, n)$}.
    \end{align*}
\end{theorem}

Two lemmas appear to be essential in Orponen's proofs of his Theorems 1.2 and 1.3. The first was the \textit{Katz-Tao Lemma} (Lemma~\ref{lemma:Katz_Tao}) claiming the existence of $(C, \delta, s)$-sets strongly covering subsets of Hausdorff dimension less than $s$. The other lemma is stated below. Roughly, it claims that any bounded set appearing at precision $\delta$ to have upper box-counting dimension no more than $s$ can be partitioned into a ``bad'' part which has small $s$-dimensional $\infty$-content, and a ``good'' part covered by some $(C, \delta, s)$-set. 

\begin{lemma}[Lemma 2.3 in \cite{Orp21}]
    \label{lemma:Lemma_2.3_of_Orponen}
    Let $0 \leq s \leq m$, $\delta > 0$, $C \geq 1$, and $X \subset \mathbb{R}^m$ be a bounded subset with $N(X, \delta) \leq C \cdot \delta^{-s}$. Then, for any $L \geq 1$, there exists a disjoint decomposition $X = X_{\operatorname{good}} \sqcup X_{\operatorname{bad}}$ such that
    \begin{enumerate}[label = (\roman*)]
        \item $\mathcal{H}^s_{\operatorname{I}}(X_{\operatorname{bad}}) \leq C_m / L$, where $C_m$ only depends on $m$, and
        \item $X_{\operatorname{good}}$ is contained in the $\delta$-neighborhood of a $(CL, \delta, s)$-set.
    \end{enumerate}
\end{lemma}

Our main goal is to prove both these lemmas using incompressibility arguments. It is already well known that Kolmogorov complexity may be used for proving combinatorial statements, including in many of the examples from Chapter 6 of \cite{LV08}. The typical advantage of an argument based on Kolmogorov complexity is in its ability to identify objects with a desired property as those which are sufficiently incompressible. And standard counting arguments show that most objects in a space should be incompressible. For simplicity, we will present these arguments over Cantor space. So let us tailor the definition of a $(C, \delta, s)$-set to this setting. 

\begin{definition}
    \label{definition:C_k_s_set}
    Fix $m \in \omega$ and let $s > 0$, $k \in \omega$, and $C \geq 1$. A finite set $P \subseteq (2^{<\omega})^m$ is called a \textit{$(C, 2^{-k}, s)$-set} if for any $\tau \in (2^{\leq k})^m$,
    \begin{align*}
        \abs{P \cap \bbrk{\tau}} \leq C \cdot \p{\frac{2^{-\operatorname{len}(\tau)}}{2^{-k}}}^s.
    \end{align*}
\end{definition}

Note that in this setting, if $P \subseteq (2^k)^m$ is a collection of $m$-tuple of strings all with equal length $k$, then $[P]$ is the $(\sqrt{m} \cdot 2^{-k})$-neighborhood about $P$ (regarding $P$ as the set of all its elements extended by an infinite tail of zeros).

We start with Lemma~\ref{lemma:Katz_Tao}. The only change we have made to its statement is a small improvement in the constant term from $Ck^2$ to $Ck$.

\begin{proposition}[Lemma~\ref{lemma:Katz_Tao}, restated]
    Let $0 < s \leq m$ and let $X \subseteq (2^{\omega})^m$ be a subset with $\dim_{\operatorname{H}} X < s$. Then there exists a constant $C \geq 1$ depending only on $m$, $s$, and $\dim_{\operatorname{H}} X$ such that: for every $k \in \omega$, there exists a $(Ck, 2^{-k}, s)$-set $P_k$ such that the sequence $\p{B_{C_m \cdot 2^{-k}}(P_k)}_{k \in \omega}$ strongly covers $X$, where $C_m \geq 1$ only depends on $m$.
    \label{proposition:Katz_Tao_restated}
\end{proposition}
\begin{proof}
    Let $B$ be a Hausdorff oracle for $X$. Fix $C_m = \sqrt{m}$ and take $\varepsilon \in \mathbb{Q}_{> 0}$ such that $\operatorname{dim}_{\operatorname{H}} X < s - \varepsilon$. Let us begin with a natural guess for such a strong cover. Start by defining the set of $m$-tuples of finite binary strings with bounded prefix complexity relative to $B$: 
    \begin{align*}
        S := \set{\sigma \in \p{2^{< \omega}}^m}{K^B(\sigma) < (s - \varepsilon) \operatorname{len}(\sigma) + K^B(\operatorname{len}(\sigma)) - c},
    \end{align*}
    where the constant $c > 0$ comes from item (ii) in Chaitin's Counting Theorem~\ref{theorem:chaitin_counting_theorem} relativized to $B$. For a fixed $j \in \omega$, denote the slice of $S$ of length $j$:
    \begin{align*}
        S_j := S \cap (2^{j})^m.
    \end{align*}
    % Let $C_m = \ceil{\sqrt{m}}$. 
    We claim that the sequence $([S_j])_{j \in \omega}$
    % which should be like doing $(B_{C_m \cdot 2^{-j}} (S_j))_j$
    must strongly cover $X$. The Point-to-Set Principle~\ref{theorem:PTS_for_H_dimension} guarantees that no $x \in X$ may have all but finitely many of its prefixes in $S$, or else $\dim_{\operatorname{H}} X \geq \dim^B(x) \geq s - \varepsilon$. Now, by the Counting Theorem~\ref{theorem:chaitin_counting_theorem}, 
    \begin{align*}
        \abs{S_j} \leq 2^{(s - \varepsilon) \cdot j}.
    \end{align*}
    Notice too that all elements in $S_j$ are mutually at a distance at least $2^{-j}$-apart, but $S_j$ is not necessarily a $(Cj, 2^{-j}, s)$-set. Instead, for each $j$, we let $S^*_j$ represent the the \textit{$s$-optimal cover} of $S_j$ in the sense of Definition~\ref{definition:optimal_cover}. That is, we find an $S^*_j$ consisting of $m$-tuples of strings of equal length at most $j$ such that $[S_j] \subseteq [S^*_j]$ and $S^*_j$ is of minimal direct $s$-weight. This minimum exists because we know the direct $s$-weight of $S_j$ satisfies:
    \begin{align*}
        \operatorname{DW}_s(S_j) = \sum_{\sigma \in S_j} 2^{- s \cdot \operatorname{len}(\sigma)} = \abs{S_j} \cdot 2^{-s \cdot j} \leq 2^{(s - \varepsilon) \cdot j} \cdot 2^{- s \cdot j} = 2^{- \varepsilon \cdot j}.
    \end{align*}
    This implies our search for $S^*_j$ should not consider strings of length shorter than $\varepsilon j / s$ (nor longer than $j$). Next, for fixed $j, k \in \omega$, stratify the elements of the optimal $s$-cover $S^*_j$ by their string-length $k$:
    \begin{align*}
        S^*_{j,k} := S^*_j \cap (2^{k})^m. 
    \end{align*}
    We first claim that each $S^*_{j,k}$ is a $(C_m, 2^{-k}, s)$-set. Indeed, for any $k \in \omega$ and $\tau \in (2^{\leq k})^m$, 
    \begin{align*}
        2^{-s \cdot k} \cdot \abs{S^*_{j,k} \cap \bbrk{\tau}} = \sum \set{2^{- s \cdot \operatorname{len}(\sigma)}}{\sigma \in S^*_{j, k} \AND \sigma \succeq \tau} \leq C_m \cdot 2^{-s \cdot \operatorname{len}(\tau)}.
    \end{align*}
    Now for each $k \in \omega$, collect into $P_k := \bigcup_{j \in \omega} S^*_{j, k}$ all the strings across the $S^*_{j, k}$. We claim that each such batch $P_k$ is a $(C k, 2^{-k}, s)$-set, where $C = C_m s / \varepsilon$. Since $S^*_{j, k} = \emptyset$ whenever $k < \varepsilon j / s$, we may more precisely write:
    \begin{align*}
        P_k = \bigcup_{j = 0}^{\floor{ks / \varepsilon}} S^*_{j, k}. 
    \end{align*}
    So, for any $\tau \in (2^{\leq k})^m$, each $S^*_{j, k}$ being a $(C_m, 2^{-k}, s)$-set implies for any $k > 0$:
    \begin{align*}
        \abs{P_k \cap \bbrk{\tau}} \leq \sum_{j = 1}^{\floor{ks / \varepsilon}} \abs{S^*_{j, k} \cap \bbrk{\tau}} \leq C_m \cdot \frac{ks}{\varepsilon} \cdot \p{\frac{2^{-\operatorname{len}(\tau)}}{2^{-k}}}^s. 
    \end{align*}
    
    Finally, we show that $X$ is indeed strongly covered by the sequence $([P_k])_{k \in \omega}$. 
    
    Fix $x \in X$. Since the $([S_j])_j$ strongly cover $X$, there exist infinitely many $j$ for which $x \in [S_j] \subseteq [S^*_j]$. For each such $j$, there exists a corresponding $k \geq \varepsilon j / s$ such that $x \in [S^*_{j, k}] \subseteq [P_k]$. Thus, it also holds that there exist infinitely many $k$ for which $x \in [P_k] $. 
\end{proof}

Now we move on to Lemma 2.3 in \cite{Orp21}. 

\begin{proposition}[Lemma 2.3 in \cite{Orp21} over Cantor Space]
    \label{proposition:effective_good_bad_decomp_Cantor}
    Let $s \in [0, m]$, $k \in \omega$, $C \geq 1$, and $X \subseteq (2^\omega)^m$ be bounded with $N(X, 2^{-k}) \leq C \cdot 2^{s \cdot k}$.
    Then for any $L \geq 1$, there exists a disjoint decomposition $X = X_{\operatorname{good}} \sqcup X_{\operatorname{bad}}$ such that
    \begin{enumerate}[label=(\roman*)]
        \item $\mathcal{H}^s_{\operatorname{I}}(X_{\operatorname{bad}}) \leq C_m/L$, where $C_m$ only depends on $m$, and
        \item $X_{\operatorname{good}}$ is contained in the $(C_m \cdot 2^{-k})$-neighborhood of a $(CL, 2^{-k}, s)$-set. 
    \end{enumerate}
\end{proposition}
\begin{proof}    
    % Let $C_m = \ceil{m^{-1/2}}$. By condition (\ref{assumption:effective_cover_bound}), the number of strings of length $k$ intersecting $X$ is bounded from above:
    % \begin{align*}
    %     \abs{\restrict{X}{k}} \leq C_m \cdot  N(X, 2^{-k}) \leq  C_m \cdot N_{\operatorname{eff}}(X, 2^{-k}) \leq C \cdot C_m \cdot 2^{ks}.
    % \end{align*}
    Let us fix a set of ``bad strings'': all those strings of length at most $k$ with an extension into $X$ and of sufficiently small complexity,
    \begin{align*}
        S_{\operatorname{bad}} := \set{\sigma \in (2^{\leq k})^m}{[\sigma] \cap X \neq \emptyset \quad \text{and} \quad K(\sigma) < s \cdot \operatorname{len}(\sigma) - \log L}.
    \end{align*}
    Take $X_{\operatorname{bad}} = [S_{\operatorname{bad}}]$ to be the open set generated by $S_{\operatorname{bad}}$.
    % Denote by $\mathcal{M}$ the subset of $S_{\operatorname{bad}}$ with all minimal length strings having none of their prefixes in $S_{\operatorname{bad}}$; this is a prefix-free set generating the same open set $[\mathcal{M}] = [S_{\operatorname{bad}}] = X_{\operatorname{bad}}$.
    Then, by the Kraft Inequality~\ref{theorem:kraft_inequality},
    \begin{align*}
        \mathcal{H}^s_{\operatorname{I}}(X_{\operatorname{bad}}) \leq \operatorname{DW}_s(S_{\operatorname{bad}}) = \sum_{\sigma \in S_{\operatorname{bad}}} 2^{- s \cdot \operatorname{len}(\sigma)} \leq \frac{1}{L} \sum_{\sigma \in S_{\operatorname{bad}}} 2^{-K(\sigma)} \leq \frac{1}{L}.
    \end{align*}
    This verifies (i) in the claim. Now, put $X_{\operatorname{good}} := X \setminus X_{\operatorname{bad}}$, and $S_{\operatorname{good}} := \restrict{X_{\operatorname{good}}}{k} = \set{\sigma \in (2^{k})^m}{[\sigma] \cap X_{\operatorname{good}} \neq \emptyset}$. Then,
    \begin{align*}
        X_{\operatorname{good}} \subseteq [\restrict{X_{\operatorname{good}}}{k}] = [S_{\operatorname{good}}],
    \end{align*}
    i.e., $X_{\operatorname{good}}$ is contained in the $(C_m \cdot 2^{-k})$-neighborhood of $S_{\operatorname{good}}$.

    It remains to show that $S_{\operatorname{good}}$ is a $(CL, 2^{-k}, s)$-set in order to conclude (ii). Fix an arbitrary $\tau \in (2^{\leq k})^m$. If $\tau \in S_{\operatorname{bad}}$, we would have $[\tau] \subseteq X_{\operatorname{bad}}$, so $S_{\operatorname{good}} \cap \bbrk{\tau} = \emptyset$. Now assume that $\tau \not\in S_{\operatorname{bad}}$. 

    By the assumption on $N(X, 2^{-k})$, we may bound the $k$-precision prefix complexity of any $x \in X$ from above via a two-part description for $\restrict{x}{k}$. Fixing an enumeration of the minimal-cardinality c.e. covering of $X$ by $2^{-k}$-balls,
    \begin{align*}
        K(\restrict{x}{k}) \leq \log N(X, 2^{-k}) + O(1) \leq s \cdot k + \log C + O(1). 
    \end{align*}
    The same bound holds for any $\sigma \in \restrict{X}{k}$:
    \begin{align}
        \label{equation:complexity_bound_for_strings_in_X_k}
        K(\sigma) \leq s \cdot \operatorname{len}(\sigma) + \log C + O(1).
    \end{align}
    Suppose $\sigma \succeq \tau$ for $\sigma \in \restrict{X}{k}$. Then by applying both the Coding Theorem~\ref{theorem:coding_theorem} and the Counting Theorem~\ref{theorem:chaitin_counting_theorem} (ii) to~\eqref{equation:complexity_bound_for_strings_in_X_k}, we deduce that $\tau$ has at most 
    \begin{align*}
        2^{K(\sigma) - K(\tau) + K(k) + O(\log \operatorname{len}(\tau)) + O(1)}
    \end{align*}
    many descriptions of length $K(\sigma) + O(\log \operatorname{len}(\tau)) + O(1)$ (each $\sigma \in \restrict{X}{k}$ extending $\tau$ computes $\tau = \restrict{\sigma}{\operatorname{len}(\tau)}$). Now, use that $\tau \not\in S_{\operatorname{bad}}$ to bound the number of extensions of $\tau$ which are ``good'':
    \begin{align*}
        \abs{S_{\operatorname{good}} \cap \bbrk{\tau}} &= \abs{\set{\sigma \in \restrict{X_{\operatorname{good}}}{k}}{\sigma \succeq \tau}}\\
        &\leq \max \set{2^{K(\sigma) - K(\tau) +  K(k) + O(\log \operatorname{len}(\tau)) + O(1)}}{\sigma \in \restrict{X}{k}}\\
        &\leq 2^{(s \cdot k + \log C + O(1)) - (s \operatorname{len}(\tau) - \log L) + o(k)}\\
        &= 2^{s(k - \operatorname{len}(\tau)) + \log C + \log L + o(k)} \\
        &\leq C' L \p{\frac{2^{-\operatorname{len}(\tau)}}{2^{-k}}}^s.
    \end{align*}
    One could adjust the definition of $S_{\operatorname{bad}}$ in the beginning to absorb the extra sub-exponential terms contributing to $C'$, resulting in a collection $S_{\operatorname{good}}$ which indeed is a $(C L, 2^{-k}, s)$-set.
\end{proof}

We discuss here how to relate the above proof with that found in \cite{Orp21}. The original proof forms $X_{\operatorname{bad}}$ as a union of all so-called \textit{heavy} dyadic cubes, i.e., those dyadic cubes $Q \in \mathcal{Q}^m$ which contain a sufficiently many smaller sub-cubes $Q'$ of a fixed side-length $2^{-k}$ intersecting with $X$:
\begin{align*}
    \bigcup \set{Q}{\abs{\set{Q'}{\ell(Q') = 2^{-k} \text{ and }  Q' \subseteq Q \text{ and }  Q' \cap X \neq \emptyset}} \geq \alpha \cdot L \cdot \p{\frac{\ell(Q)}{2^{-k}}}^s},
\end{align*}
where $\alpha$ is $C$ times a small constant depending on $m$, and $\ell(Q)$ denotes the side-length of $Q$. Following a strict translation between Euclidean and Cantor spaces (and omitting relativization to Hausdorff oracles), our proof could have defined $X_{\operatorname{bad}}$ to be the union of all infinite extensions of \textit{heavy} strings, i.e., finite strings which may be extended to sufficiently many strings of length $k$ whose infinite extension intersects with $X$:
\begin{align*}
    \brk{\set{\sigma}{\abs{\set{\tau \succeq \sigma}{[\tau] \cap X \neq \emptyset \text{ and } K(\tau) \leq s \cdot \operatorname{len}(\tau)}} \geq \alpha \cdot L \cdot \p{\frac{2^{-\operatorname{len}(\sigma)}}{2^{-k}}}^s}}.
\end{align*}
The only liberty we have taken in the above translation is to add an incompressibility condition on $\tau$, making the ``bad'' set even more restrictive. 

In particular, any such heavy string has many \textit{long descriptions} using each of the length $k$ strings $\tau$ of bounded complexity extending $\sigma$ and intersecting with $X$. Together with $\operatorname{len}(\sigma)$, such a string $\tau$ suffices to compute $\sigma$ by truncation. The Coding Theorem~\ref{theorem:coding_theorem} guarantees that $\sigma$ having many such long descriptions implies that $\sigma$ has at least one short description. In particular, $\sigma$ having at least $\alpha \cdot L \cdot 2^{s(k - \operatorname{len}(\sigma)) + K(k)}$ many descriptions $\langle\tau^*, \operatorname{len}(\sigma) \rangle$ of length at most $K(\tau) + K(\operatorname{len}(\sigma)) + O(1) \leq sk + o(k)$ guarantees 
\begin{align*}
    K(\sigma) \leq \p{s \cdot k + o(k)} - \log \brk{\alpha \cdot L \cdot 2^{s(k - \operatorname{len}(\sigma)) + K(k)}} = s \cdot \operatorname{len}(\sigma) - \log L - \log \alpha + o(k).
\end{align*}
This essentially places $\sigma$ in the canonical $S_{\operatorname{bad}}$ set from our proof, or its open cylinder $\brk{\sigma}$ in the canonical $X_{\operatorname{bad}}$ from our proof. That is, we may view the effective analog of $S_{\operatorname{bad}}$ as defined by Orponen as a subset of our bad set of strings.

Finally, this result effectivizes for a restricted class of subsets. For instance,

\begin{corollary}
    \label{corollary:effective_good_bad_decomp_Cantor}
    Let $s \in [0, m] \cap \mathbb{Q}$, $k \in \omega$, $C \geq 1$, and $X \subseteq (2^\omega)^m$ be a bounded $\Sigma^0_1$-class with $N(X, 2^{-k}) \leq C \cdot 2^{s \cdot k}$.
    Then for any rational $L \geq 1$, there exists a disjoint decomposition $X = X_{\operatorname{good}} \sqcup X_{\operatorname{bad}}$ with $X_{\operatorname{good}} \in \Pi^0_1$ and $X_{\operatorname{bad}} \in \Sigma^0_1$ such that
    \begin{enumerate}[label=(\roman*)]
        \item $\mathcal{H}^s_{\operatorname{I}}(X_{\operatorname{bad}}) \leq C_m/L$, where $C_m$ only depends on $m$, and
        \item $X_{\operatorname{good}}$ is contained in the $(C_m \cdot 2^{-k})$-neighborhood of a $(CL, 2^{-k}, s)$-set. 
    \end{enumerate}
\end{corollary}

% We may conclude the general result from its effectivization.

% \begin{corollary}[C.f. Lemma 2.3 in \cite{Orp21}]
%     Let $0 \leq s \leq m$, $k \in \omega$, $C \geq 1$, and $X \subseteq (2^\omega)^m$ be a bounded set with 
%     \begin{align*}
%         \label{assumption:cover_bound}
%         N(X, 2^{-k}) \leq C \cdot 2^{s \cdot k}.
%     \end{align*}
%     Then for any $L \geq 1$, there exists a disjoint decomposition $X = X_{\operatorname{good}} \sqcup X_{\operatorname{bad}}$ such that
%     \begin{enumerate}[label=(\roman*)]
%         \item $\mathcal{H}^s_{\operatorname{I}}(X_{\operatorname{bad}}) \leq 1/L$, and
%         \item $X_{\operatorname{good}}$ is contained in the $2^{-k}$-neighborhood of a $(CL, 2^{-k}, s)$-set. 
%     \end{enumerate}
%     \label{cor:good_bad_decomp}
% \end{corollary}
% \begin{proof}
%     By approximation by closed supersets, it suffices to consider $X$ to be a $\Delta^0_1$-class relative to some oracle $B \in 2^{\leq \omega}$. Indeed, for any $\varepsilon > 0$ and closed superset $\widetilde{X} \supseteq X$ satisfying $\dim_{\operatorname{H}} \widetilde{X} = \dim_{\operatorname{H}} X + \varepsilon$, we have that $\widetilde{X}$ is effectively closed with respect to some oracle, hence computable in the Turing-jump of that oracle. Moreover, we may ask that this oracle also encode some cover witnessing $N(\widetilde{X}, 2^{-k})$.   
    
%     Thus, we may apply Proposition~\ref{proposition:effective_good_bad_decomp_Cantor} relativized to $B$ to decompose $\widet$ the decomposition $\widetilde{X} = \widetilde{X}_{\operatorname{good}} \sqcup \widetilde{X}_{\operatorname{bad}}$. Passing to only the parts intersecting with $X$, decompose $X$ as follows:
%     \begin{align*}
%         X = X_{\operatorname{good}} \sqcup X_{\operatorname{bad}} := (\widetilde{X}_{\operatorname{good}} \cap X) \sqcup (\widetilde{X}_{\operatorname{bad}} \cap X). 
%     \end{align*}
%     Still, $X_{\operatorname{bad}}$ is of small $s$-dimensional measure, and $X_{\operatorname{good}}$ is contained in the $2^{-k}$-neighborhood of a $(CL, 2^{-k}, s)$-set.
%     {\color{red} If we are to approximate by closed supersets, then I don't konw how to still guarantee that the $N_eff$ bound holds for that superset... so can't apply the effective theorem?}
% \end{proof}

\section{Projective and Pinned-Distance Sets} \label{section:Projective_and_Pinned_Distance_Sets}

In this section we present a simple effective argument which implies results for Hausdorff dimension in multiple geometric setups. These geometries will include pinned-distance sets, radial projections, and orthogonal projections. 

\begin{definition}
    Fix $0 \leq n \leq m$. Given any norm $\norm[*]{\cdot}$ on $\mathbb{R}^m$, subset $X \subseteq \mathbb{R}^m$, point $z \in \mathbb{R}^m$, and subspace $V \in G(m, n)$, denote
    \begin{itemize}
        \item The \textit{pinned-distance} of $X$ about $z$ by $\Delta^*_z(X) := \set{\norm[*]{x - z}}{x \in X}$;
        \item The \textit{radial projection} of $X$ about $z$ by $P_z^*(X) := \set{\frac{x - z}{\norm[*]{x - z}}}{x \in X}$;
        \item The \textit{orthogonal projection} of $X$ onto $V$ by $\Pi^*_V(X) := \set{\norm[*]{\operatorname{proj}_V(x)}}{x \in X}$;
    \end{itemize}
    all with respect to the norm $\norm[*]{\cdot}$. Simply write $\Delta_z(X)$, $P_z(X)$, and $\Pi_V(X)$ when taken with respect to the standard $L^2$-norm. 
\end{definition}

We start by recalling a result about pinned-distance sets by D.~Stull \cite{Stu22_pinned}. 
\begin{theorem}[Theorem 1 in \cite{Stu22_pinned}] 
    \label{theorem:Stull_pinned_distance_theorem}
    Let $X \subseteq \mathbb{R}^2$ be an analytic set with Hausdorff dimension strictly greater than one. Then, for all $z \in \mathbb{R}^2$ outside a set of Hausdorff dimension at most one, 
    \begin{align*}
        \dim_{\operatorname{H}}\p{\Delta_z X} \geq \frac{\dim_{\operatorname{H}}X}{4} + \frac{1}{2}.
    \end{align*}
\end{theorem}

In 2023, I.~Altaf, R.~Bushling, and B.~Wilson proved a related lower-bound for the Hausdorff dimension of pinned-distance sets \cite{ABW23}.

\begin{theorem}[Theorem 1.2 in \cite{ABW23}] \label{theorem:Bushling_pinned_distance_theorem}
    For each $m \in \omega$, norm $\norm[*]{\cdot}$ on $\mathbb{R}^m$, point $z \in \mathbb{R}^m$, and subset $X \subseteq \mathbb{R}^m$,
    \begin{align*}
        \dim_{\operatorname{H}}\p{\Delta^*_z X} \geq \dim_{\operatorname{H}}X - (m - 1).
    \end{align*}
\end{theorem}

It is easy to see that under the conditions of Theorem~\ref{theorem:Stull_pinned_distance_theorem}, D. Stull's lower bound is stronger than that of Theorem~\ref{theorem:Bushling_pinned_distance_theorem}. But the latter applies more generally: there are no exceptional directions, no dimensionality assumptions on $X$, nor any specificity to the $L^2$-norm.

Recall that the Marstrand-Mattila Projection Theorem~\ref{theorem:marstrand_projection_original} bounds from below the Hausdorff dimension of the orthogonal projection of any analytic set $X \subseteq \mathbb{R}^m$ along all but a $\gamma_{m, n}$-null set of subspaces from $G(m, n)$. Is there a lower bound for orthogonal projections which, as in Theorem~\ref{theorem:Bushling_pinned_distance_theorem}, applies without any exceptional directions?

To answer this question, we recall a well-known bound for the Hausdorff dimension of a set which can be decomposed into two components, and provide an effective proof based on that of Theorem~\ref{theorem:Bushling_pinned_distance_theorem}.  

\begin{theorem}
    Let $m, n, \ell\in \omega$ and $X \subseteq \mathbb{R}^m$, $U \subseteq \mathbb{R}^n$, and $W \subseteq \R^\ell$. Suppose there is a locally Lipschitz continuous map $f: U \times W \to \mathbb{R}^m$ surjective on $X$. Then
    \begin{align*}
        \dim_{\operatorname{H}} X \leq \dim_{\operatorname{H}} U + \dim_{\operatorname{P}} W.
    \end{align*}
    \label{theorem:meta_theorem_two_part_sets_dimension_bound}
\end{theorem}
Classically, Theorem~\ref{theorem:meta_theorem_two_part_sets_dimension_bound} is proved by the following sequence of inequalities:
\begin{align*}
    \dim_{\operatorname{H}} X &\leq \dim_{\operatorname{H}} f(U \times W) \leq \dim_{\operatorname{H}}(U \times W) \leq \dim_{\operatorname{H}} U + \dim_{\operatorname{P}} W.
\end{align*}
The first inequality follows from monotonicity of Hausdorff dimension under inclusion. The second inequality follows from monotonicity of Hausdorff dimension under locally Lipschitz continuous mappings (see Corollary 2.4 in \cite{Fal90}). The final inequality originally follows from Theorem 3 in \cite{Tri82}, which takes the most work to show.

With respect to computability, we may make use of the fact that any (locally) Lipschitz continuous map is computable in some oracle.

\begin{proof}
    Let $B \geq_{\operatorname{T}} f$ be a Hausdorff oracle for both $X$ and $U$, as well as a packing oracle for $W$. Notice that for any point $x \in X$ there exists $u \in U$ and $w \in W$ such that for any precision-level $r \in \omega$,
    \begin{align*}
        K^B_r(x) = K^B_r(f(u, w)) \leq K^B_r(u, w) + o(r) \leq K^B_r(u) + K^B_r(w) + o(r).
    \end{align*}
    Let $\varepsilon > 0$. Since $B$ is a Hausdorff oracle for $X$, there exists $x \in X$ such that $\dim_{\operatorname{H}} X \leq \dim^B(x) + \varepsilon$. Then, by the choice of $B$ and the subadditivity of $K$:
    \begin{align*}
        \dim_{\operatorname{H}} X - \varepsilon
        &\leq \dim^B(x)\\
        &= \liminf_{r \to \infty} \frac{K^B_r(x)}{r}\\
        &\leq \liminf_{r \to \infty} \frac{K^B_r(u) + K^B_r(w)}{r}\\
        &\leq \dim^B(u) + \Dim^B(w)\\
        &\leq \dim_{\operatorname{H}} U + \dim_{\operatorname{P}} W.
    \end{align*}
\end{proof}

Let us use Theorem~\ref{theorem:meta_theorem_two_part_sets_dimension_bound} to recreate the proof of Theorem~\ref{theorem:Bushling_pinned_distance_theorem}.

\begin{proof}[Proof of Theorem~\ref{theorem:Bushling_pinned_distance_theorem}]
    Take $U = \Delta^*_z(X)$ and $W = B^*_1(0)$ the unit ball in $\mathbb{R}^m$ with respect to the norm $\norm[*]{\cdot}$. Notice that for any $X \ni x \neq z$, we may express $x$ as follows:
    \begin{align*}
        x = \norm[*]{x - z} \frac{x - z}{\norm[*]{x - z}} + z.
    \end{align*}
    Since $\norm[*]{x - z} \in \Delta^*_z(X)$ and $\frac{x - z}{\norm[*]{x - z}} \in B^*_1(0)$, we have found a mapping $f : U \times W \to \mathbb{R}^m$, $f: (u, w) \mapsto u w + z$ which surjects onto $X$, is locally Lipschitz continuous, and is computable in $\seq{z}$ (if $x = z$, simply return $z$). Noting that $\dim_{\operatorname{H}} W = \dim_{\operatorname{P}} W = m - 1$, Theorem~\ref{theorem:meta_theorem_two_part_sets_dimension_bound} gives:
    \begin{align*}
        \dim_{\operatorname{H}} X \leq \dim_{\operatorname{H}}(\Delta^*_z(X)) + (m - 1).
    \end{align*}
\end{proof}

\begin{proposition}
    \label{proposition:Bushling_type_projection_theorem}
    For each $m \in \omega$, norm $\norm[*]{\cdot}$ on $\mathbb{R}^m$, direction $e \in \mathbb{S}^{m - 1}$, and subset $X \subseteq \mathbb{R}^m$,
    \begin{align*}
        \dim_{\operatorname{H}}\p{\Pi^*_e X} \geq \dim_{\operatorname{H}} X - (m - 1),
    \end{align*}
    where $\Pi^*_e = \Pi^*_{\operatorname{span}(e)}$.
\end{proposition}

\begin{proof}
    Take $U = \Pi^*_e(X) \subseteq \mathbb{R}$ and $W = \mathcal{O}_e := \set{w \in \mathbb{R}^m}{w \perp e} \in G(m, m - 1)$ the orthogonal subspace to $e$. Notice that for any $x \in X$, we may express $x$ as follows:
    \begin{align*}
        x = \norm[*]{\operatorname{proj}_e(x)} e + (x - \norm[*]{\operatorname{proj}_e(x)} e).
    \end{align*}
    Since $\norm[*]{\operatorname{proj}_e(x)} \in \Pi^*_e(X)$ and $x - \norm[*]{\operatorname{proj}_e(x)} e \in \mathcal{O}_e$, the mapping $f : U \times W \to \mathbb{R}^m$, $f : (u, w) \mapsto u e + w$ surjects onto $X$, is Lipschitz continuous, and is computable in $\seq{e}$. Furthermore, $\dim_{\operatorname{H}} W = \dim_{\operatorname{P}} W = m - 1$. So, by Theorem~\ref{theorem:meta_theorem_two_part_sets_dimension_bound},
    \begin{align*}
        \dim_{\operatorname{H}} X \leq \dim_{\operatorname{H}}(\Pi^*_e(X)) + (m - 1).
    \end{align*}
\end{proof}

Since there is a bi-Lipschitz continuous map $\norm[*]{\cdot}$ translating between $\Pi^*_e(X)$ and $\operatorname{proj}_e(X)$, Proposition~\ref{proposition:Bushling_type_projection_theorem} also gives
\begin{align*}
    \dim_{\operatorname{H}}\p{\operatorname{proj}_e(X)} \geq \dim_{\operatorname{H}}X - (m - 1).
\end{align*}

We obtain another result of this kind for radial projections with a well-known corollary. We refer \cite{Orp19} to as a standard source.

\begin{proposition}
    \label{proposition:decomposition_result_for_radial_projections}
    For each $m \in \omega$, norm $\norm[*]{\cdot}$ on $\mathbb{R}^m$, point $z \in \mathbb{R}^m$, and subset $X \subseteq \mathbb{R}^m$,
    \begin{align*}
        \dim_{\operatorname{H}}\p{P^*_z X} \geq \dim_{\operatorname{H}} X - 1.
    \end{align*}
\end{proposition}
The proof relies on the same geometry as that of Theorem~\ref{theorem:Bushling_pinned_distance_theorem}, but with distinct interpretations for the sets involved in the decomposition.

\begin{proof}
    Take $U = P^*_z(X)$ and $W = \mathbb{R}$. Notice that for any $X \ni x \neq z$, we may express $x$ as follows:
    \begin{align*}
        x = \frac{x - z}{\norm[*]{x - z}} \cdot \norm[*]{x - z} + z.
    \end{align*}
    Since $\frac{x - z}{\norm[*]{x - z}} \in P^*_z(X)$ and $\norm[*]{x - z} \in \mathbb{R}$, we have found a mapping $f : U \times W \to \mathbb{R}^m$, $f: (u, w) \mapsto u w + z$ which surjects onto $X$, is locally Lipschitz continuous, and is computable in $\seq{z}$ (if $x = z$, simply return $z$). Noting that $\dim_{\operatorname{H}} W = \dim_{\operatorname{P}} W = 1$, Theorem~\ref{theorem:meta_theorem_two_part_sets_dimension_bound} gives:
    \begin{align*}
        \dim_{\operatorname{H}} X \leq \dim_{\operatorname{H}}(P^*_z(X)) + 1.
    \end{align*}
\end{proof}

\begin{definition}
    Let $z \in \mathbb{R}^m$. A subset $X \subseteq \mathbb{R}^m$ is said to be \textit{totally invisible from $z$} if $\dim_{\operatorname{H}}(P_z(X \setminus \seq{z})) = 0$. Let $\operatorname{Inv}_{\operatorname{T}}(X)$ be the collection of all points from which $X$ is totally invisible. 
\end{definition}

Suppose a subset $X$ of Hausdorff dimension greater than one were to be totally invisible from some point $z$; then Proposition~\ref{proposition:decomposition_result_for_radial_projections} would imply $\dim_{\operatorname{H}}(P_z(X \setminus \seq{z})) \geq \dim_{\operatorname{H}} X - 1$, which is positive, a contradiction. This gives the well-known corollary.

\begin{corollary}
    For any subset $X \subseteq \mathbb{R}^m$ with $\dim_{\operatorname{H}} X > 1$, we have $\operatorname{Inv}_{\operatorname{T}}(X) = \emptyset$.
\end{corollary}

In summary, Theorem~\ref{theorem:meta_theorem_two_part_sets_dimension_bound} may be applied to various geometries. Theorem~\ref{theorem:Bushling_pinned_distance_theorem} is an application to pinned-distance sets and the unit ball. Dually, Proposition~\ref{proposition:decomposition_result_for_radial_projections} is an application to radial projections and the real line. Likewise, Proposition~\ref{proposition:Bushling_type_projection_theorem} is an application to orthogonal projections and their orthogonal hyperplanes.

\section{Subsets of Exact Measure}
\label{section:subsets_of_exact_measure}

A standard result in geometric measure theory is one by A.~Besicovitch regarding for each closed set in Euclidean space of positive $\mathcal{H}^s$-measure finding a subset whose $\mathcal{H}^s$-measure is non-zero and finite \cite{Bes52}. The result was quickly shown to hold more generally for analytic (or Souslin) sets by R.~Davies \cite{Dav52}. Such a result is trivial when $s$ equals the full dimension of the ambient space, since then $\mathcal{H}^s$ is a constant multiple of Lebesgue measure on that space. Without such a correspondence for smaller values of $s$, one must get around the fact that not all sets of infinite $\mathcal{H}^s$-measure are $\sigma$-finite for $\mathcal{H}^s$ (i.e., are not the countable union of subsets of finite $\mathcal{H}^s$-measure). Besicovitch circumvents this by using a net measure constructed over the net of dyadic cubes $\mathcal{Q}^m$ over Euclidean space $\mathbb{R}^m$ which is comparable to $\mathcal{H}^s$. As discussed in Chapter~\ref{chapter:AIT_on_Net_Spaces}, D.~Larman extended nets and net measures to generic metric spaces \cite{Lar67}. Soon after, C.~Rogers and Davies managed to extend Besicovitch's result on the existence of finite measure subsets to compact subsets of a net space under some weak assumptions about the net and net measure \cite{Rog70}, presented below:

\begin{theorem}[Theorem 54 of \cite{Rog70}]
    \label{theorem:besicovitch_by_rogers_for_nets}
    Let $\mathcal{N}$ be a Rogers net on a metric space $(\Omega, d)$, and $\rho$ be a net premeasure for $\mathcal{N}$. Suppose further that
    \begin{enumerate}[label = (\roman*)]
        \item $\rho$ is finite-valued on $\mathcal{N}$,
        \item Each $x \in \Omega$ is $\mathcal{H}^\rho$-finite, and
        \item If $(F_n)_{n \in \infty}$ is a decreasing sequence of compact subsets of $\Omega$ with $\mathcal{H}^\rho\p{\bigcup_n F_n} = 0$, it holds for each $\delta > 0$ that $\mathcal{H}^\rho_\delta(F_n) \to 0$ as $n \to \infty$.
    \end{enumerate}
    Then, any compact set $F \subseteq \Omega$ with positive $\mathcal{H}^\rho$-measure contains a subset of positive, finite $\mathcal{H}^\rho$-measure.
\end{theorem}

Rogers and Davies used their Increasing Sets Lemma~\ref{theorem:increasing_sets_lemma} and stability result~\ref{theorem:stability_for_net_measures} to prove Theorem~\ref{theorem:besicovitch_by_rogers_for_nets}. 

The goal of the present section is to obtain a stronger result for a slightly more restricted class of nets and net measures. In particular, we will prove a density result which applies to any finitely-branching, layered-disjoint net with consistently shrinking diameters across the rank levels and which generates the topology of the full space. And this result will apply for any Hausdorff net premeasure (i.e., net premeasure induced by a dimension function in the sense of Section~\ref{section:Measures_and_Hausdorff_Dimension}). We originally obtained the result over Cantor space by translating Besicovitch's original argument to the tree-representations of $\boldsymbol{\Pi}^0_1$-classes. This proof method then generalized to net spaces with the properties stated above. 

\begin{definition}
    Call a net $\mathcal{N}$ on a metric space $(\Omega, d)$ \textit{scaling} if there exists a monotonically-decreasing function $\mathsf{d} : \omega \to (0, +\infty]$ with $\mathsf{d}(r) \to 0$ as $r \to \infty$ such that $\operatorname{rk}(N) = r$ implies $\operatorname{diam}(N) = \mathsf{d}(r)$. 
\end{definition}

Many standard examples of nets are scaling, including the net of dyadic cubes $\mathcal{Q}^m$ on $\mathbb{R}^m$ with $\mathsf{d}(r) = \sqrt{m} \cdot 2^{-r}$, or the net of cylinder sets $\mathcal{B}$ on $2^\omega$ with $\mathsf{d}(r) = 2^{-r}$. Notice as well that the Rogers net guaranteed in Theorem~\ref{thm:Larman_net_on_finite_dimensional_compact_metric} for a compact metric space is also scaling for $\mathsf{d}(r) = 2^{-r}$. 

\begin{definition}
    For a fixed $0 < b \in \omega$, a net $\mathcal{N}$ on a metric space $(\Omega, d)$ is \textit{$b$-branching} if $\mathcal{N}$ has a single root and each net element has exactly $b$-many children. 
\end{definition}

Any $b$-branching net admits an indexing $\mathcal{N} = \set{N_\sigma}{\sigma \in b^{<\omega}}$ respecting the containment relation, following the addresses in the associated graph $\mathcal{G}$ as in Proposition~\ref{prop:tree_characterization_of_net}. 

A collection of subsets of a topological space $\Omega$ is said to be a \textit{subbasis of $\Omega$} if it generates the same topology on $\Omega$. So, a net $\mathcal{N}$ is a subbasis of $(\Omega, d)$ whenever the topology it generates on $\Omega$ is compatible with the topology induced by the metric $d$. 

Let $\mathcal{N}$ be a scaling, $b$-branching, layered-disjoint net on $(\Omega, d)$ with Hausdorff net premeasure $\rho_h$ based on a dimension function $h$. Suppose $F \subseteq \Omega$ is a non-empty, compact class.

As $F$ is compact and $\mathcal{N}$ is $b$-branching, the net axioms (\textbf{M3}) and (\textbf{N4}$'$) guarantee that we may represent $F$ as the set of paths through a tree $T_F \subseteq b^{<\omega}$. That is, it holds that any $x \in F$ may be uniquely associated with the path-wise intersection $\bigcap_{\sigma \in p} N_\sigma$, where $p \in [T_F]$ and $N_\sigma \in \mathcal{N}$ is the net element at address $\sigma \in b^{<\omega}$ in the associated graph $\mathcal{G}$.

Using $T_F$, we may represent the set of all non-empty, compact subclasses of $F$ as its own $\boldsymbol{\Pi}^0_1$-class $\mathcal{P}_F \subseteq 2^{\omega}$. To do this, we define the tree $T_{\mathcal{P}} \subseteq 2^{<\omega}$ coding $\mathcal{P}_F$ as follows. Identify $2^{<\omega}$ with $\omega$ using a standard computable bijection with the property that for each $n \in \omega$, all strings in $b^{<\omega}$ of length $n$ are enumerated before those of length $n + 1$. For any $\zeta \in 2^{<\omega}$, define $\zeta \in T_{\mathcal{P}}$ if $\zeta$ satisfies the following conditions for all $\sigma \in b^{<\omega}$ with $\zeta(\sigma) = 1$:
\begin{enumerate}[label = (\roman*)]
    \item $\zeta(\tau) = 1$ for all $\tau \preceq \sigma$,
    \item $\sigma \in T_F$,
    \item If $\zeta$ is defined on all immediate successors $\sigma^\smallfrown i$ where $i \in \seq{0, ..., b - 1}$, then $\zeta(\sigma^\smallfrown i) = 1$ for at least one choice of $i$.
    % this last item should be modified if doing "at most b branching" appropriately
\end{enumerate}

It is straightforward to verify that these conditions make $T_{\mathcal{P}}$ a tree. Take any infinite path $Z \in [T_{\mathcal{P}}]$ through $T_{\mathcal{P}}$, and associate to it the set $T_Z = \set{\sigma \in b^{<\omega}}{Z(\sigma) = 1}$. By condition (i) above, $T_Z$ is a tree. By condition (ii), $T_Z$ is a sub-tree of $T_F$. And by condition (iii), $T_Z$ is pruned (i.e., each elements of $T_Z$ has a proper extension in $T_Z$). Further associate to $Z$ the compact set $E_Z \subseteq \Omega$ defined by
\begin{align*}
    x \in E_Z \iff (\exists p \in [T_Z]) \brk{\bigcap_{\sigma \in p} N_\sigma = \seq{x}}.
\end{align*}
For any $Z \in [T_{\mathcal{P}}]$, since $T_Z$ is a pruned sub-tree of $T_F$ and $F$ is compact, the associated set $E_Z$ is a non-empty, compact subset of $F$. And any non-empty, compact subset of $F$ may be represented by an infinite, pruned sub-tree of $T_F$, so we have our desired correspondence: $\mathcal{P}_F = [T_{\mathcal{P}}]$. 

This representation helps to prove the following result.

\begin{theorem}
    \label{theorem:existence_of_compact_subset_exact_for_premeasure}
    Let $\mathcal{N}$ be a scaling, $b$-branching, layered-disjoint net and subbasis on $(\Omega, d)$ and having net Hausdorff premeasure $\rho$. Suppose $F \subseteq \Omega$ is compact with $(\restrict{\mathcal{H}}{\mathcal{N}})^\rho(F) = +\infty$. Then there exists a compact subset $E \subseteq F$ of arbitrarily large, finite $(\restrict{\mathcal{H}}{\mathcal{N}})^\rho$-measure.
\end{theorem}

Any class $F$ with $(\restrict{\mathcal{H}}{\mathcal{N}})^\rho(F) > 0$ is certainly non-empty, so we may construct the associated class $\mathcal{P}_F$ as described above. Since $2^{\omega}$ is a complete metric space and $\mathcal{P}_F \subseteq 2^{\omega}$ is closed, $\mathcal{P}_F$ may be viewed as a complete metric space in its own right.

Our general outline will now be as follows:
\begin{enumerate}[label = (\arabic*)]
    \item Find a non-empty, compact subset $\mathcal{S} \subseteq \mathcal{P}_F$ such that if $Z \in \mathcal{S},$ then the associated subclass $E_Z \subseteq F$ satisfies $(\restrict{\mathcal{H}}{\mathcal{N}})^\rho(E_Z) > 0$.
    \item Find a countable sequence of non-empty, open sets $(\mathcal{U}^c_n)_n$ in $\mathcal{P}_F$ such that if $Z \in \bigcap_{n \in \omega} \mathcal{U}^c_n$, then the associated subclass $E_Z \subseteq F$ satisfies $(\restrict{\mathcal{H}}{\mathcal{N}})^\rho(E_Z) < +\infty$.
    \item Show that each $\mathcal{U}^c_n$ is dense in the induced subspace $\mathcal{S} \subseteq \mathcal{P}_F$.
\end{enumerate}

Since $\mathcal{P}_F$ is complete, the closed subspace $\mathcal{S}$ in (1) will also be complete. And if each $\mathcal{U}^c_n$ is indeed dense in $\mathcal{S}$, then by the Baire Category Theorem, the intersection $\bigcap_{n \in \omega}\mathcal{U}^c_n$ will be dense in $\mathcal{S}$ as well. In particular, $\mathcal{S} \cap \bigcap_{n \in \omega}\mathcal{U}^c_n$ will be non-empty, and for any $Z$ in this intersection, the associated subclass $E_Z \subseteq F$ will satisfy $0 < (\restrict{\mathcal{H}}{\mathcal{N}})^\rho(E_Z) < +\infty$, as desired.

Suppose $U \subseteq b^{<\omega}$. Understand $\operatorname{DW}_\rho(U)$ to evaluate the direct $\rho$-weight of the collection of net elements $\set{N_\sigma}{\sigma \in U}$ whose addresses are in $U$.

For any natural numbers $n \leq m$, define a function $\phi_{n, m}: \mathcal{P}_F \to \mathbb{R}$ as follows:
\begin{align*}
    \phi_{n, m}(Z) = \min_{U \subseteq b^{<\omega}} \set{\operatorname{DW}_\rho(U)}{(\forall \sigma \in U)\brk{n \leq \operatorname{len}(\sigma) \leq m} \text{ and } (\forall \sigma \in T_Z)(\exists \tau \in U)[\tau \preceq \sigma]}.
\end{align*}
That is, $\phi_{n, m}(Z)$ is defined as the minimum direct $\rho$-weight of any cover $U \subseteq b^{<\omega}$ of $[T_Z]$ by strings of length between $n$ and $m$, inclusive. There are only finitely many $U$ considered in the computation of $\phi_{n, m}(Z)$. Since $\mathcal{N}$ is scaling, any such sequence $U$ covering $[T_Z]$ may also be viewed as giving the minimal direct $\rho$-weight cover of $E_Z$ by net elements having diameter between $\mathsf{d}(m)$ and $\mathsf{d}(n)$, inclusive.

Fixing $c > 0$, define for each $n \in \omega$ the set
\begin{align*}
    \mathcal{U}^c_n := \set{Z \in \mathcal{P}_F}{(\exists m \geq n)[\phi_{n, m}(Z) < c]}.
\end{align*}
For each $n \in \omega$, it only takes finitely bits of $Z$ to determine whether $\phi_{n, m}(Z) < c$, so each $\set{Z}{\phi_{n, m}(Z) < c}$ is clopen in the product topology of $2^\omega$. As the union of all such sets for $m \geq n$ intersected with $\mathcal{P}_F$, the set $U^c_n$ is open in $\mathcal{P}_F$. Moreover, each $U^c_n$ is non-empty, for if $Z \in \mathcal{P}_F$ codes the singleton set $E_Z = \seq{x}$, then since $\mathcal{N}$ is scaling, we have $\phi_{n, m}(Z) = h(\mathsf{d}(m))$, as witnessed by the unique length-$m$ string $\sigma$ with $Z(\sigma) = 1$. So, pick $m \gg n$ large enough so that $h(\mathsf{d}(m)) < c$ to witness that $Z \in \mathcal{U}^c_n$. 

Finally, we observe that if $Z \in \bigcap_{n \in \omega}\mathcal{U}^c_n$ with $c > 0$ being fixed, then $(\restrict{\mathcal{H}}{\mathcal{N}})^\rho(E_Z) \leq c < +\infty$. Indeed, take any $\delta > 0$, and consider $n \in \omega$ for which $\mathsf{d}(n) < \delta$. By definition of $Z \in \mathcal{U}^c_n$, there is some $m \geq n$ for which $\phi_{n, m}(Z) < c$, implying there exists a $\mathsf{d}(n)$-cover of $E_Z$ using only elements of $\mathcal{N}$ and of direct $\rho$-weight less than $c$. Thus $(\restrict{\mathcal{H}}{\mathcal{N}})^\rho_\delta(E_Z) < c$. Taking the limit as $\delta \to 0$ gives $(\restrict{\mathcal{H}}{\mathcal{N}})^\rho(E_Z) \leq c$.

We now define the set $\mathcal{S}$ to witness being of positive $(\restrict{\mathcal{H}}{\mathcal{N}})^\rho$-measure. 
% Here, we use both compactness and the fact that the elements of $\mathcal{P}_F$ only represent pruned trees. 
For each $c > 0$ and $n \in \omega$, define the set:
\begin{align*}
    \mathcal{S}^c_n := \mathcal{P}_F \setminus \mathcal{U}^c_n = \set{Z \in \mathcal{P}_F}{(\forall m \geq n) [\phi_{n,m}(Z) \geq c]}.
\end{align*}
Note that each $\mathcal{S}^c_n$ is closed in $\mathcal{P}_F$ as the complement of $U^c_n$ in $\mathcal{P}_F$. 

\begin{lemma}
    \label{lemma:Hausdorff_measure_characterization_of_S_c_n}
    For any $Z \in \mathcal{P}_F$, $c > 0$, and $n \in \omega$,
    \begin{align*}
        (\restrict{\mathcal{H}}{\mathcal{N}})^\rho_{\mathsf{d}(n)}(E_Z) \geq c \iff Z \in \mathcal{S}^c_n.
    \end{align*}
\end{lemma}
\begin{proof}
    Suppose $(\restrict{\mathcal{H}}{\mathcal{N}})^\rho_{\mathsf{d}(n)}(E_Z) \geq c$. Then, any $\mathsf{d}(n)$-cover of $E_Z$ by net elements has direct $\rho$-weight at least $c$. By definition and the assumption that $\mathcal{N}$ is scaling, this includes all covers by net elements with addresses of length at least $n$. So, for any $m \geq n$, we have $\phi_{n, m}(Z) \geq c$, hence $Z \in S^c_n$. 

    In the reverse direction, assume $Z \in S^c_n$. Consider an arbitrary $\mathsf{d}(n)$-cover $U$ of $E_Z$ by net elements. Since $E_Z$ is compact, it follows by $\mathcal{N}$ being a subbasis of $(\Omega, d)$ that
    % Alexander's Subbasis Theorem 
    this cover has a finite subcover $U' = \seq{\sigma_1, ..., \sigma_t }$. 
    % {\color{red} I might be able to get rid of this subbasis assumption by mimicking a step of $K \mapsto K^*$ in Rogers' Theorem 54 and appealing to the Increasing Sets Lemma.} 
    Since $\mathcal{N}$ is scaling, each $\sigma_i \in b^{<\omega}$ is of length $\operatorname{len}(\sigma_i) \geq n$. Call $m$ the maximum length of all such $\sigma_i$. Without loss of generality, $U'$ may be restricted to just those $\sigma_i$ with an extension in $T_Z$, for any $\sigma_i$ without this property would not be needed in the cover of $E_Z$ and could thus be deleted from $U'$. 
    
    We claim that $U'$ is one of the covers considered in computing $\phi_{n, m}(Z)$. Indeed, since $Z$ codes a pruned tree, any $\sigma \in b^m$ with $Z(\sigma) = 1$ must have some extension in $E_Z$, which in turn must have some prefix $\sigma_i \in U'$ as $U'$ is a cover of $E_Z$ containing strings of length at most $m$. 
    
    So, by the assumption that $\phi_{n, m}(Z) \geq c$, it follows that $\operatorname{DW}_\rho(U') \geq c$. And, by monotonicity, the original cover $U$ of $E_Z$ must also have $\operatorname{DW}_\rho(U) \geq c$. Since this $\mathsf{d}(n)$-cover by net elements was arbitrary, we may conclude $(\restrict{\mathcal{H}}{\mathcal{N}})^\rho_{\mathsf{d}(n)}(E_Z) \geq c$, as claimed.
\end{proof}
    
By the assumption $(\restrict{\mathcal{H}}{\mathcal{N}})^\rho(F) = +\infty$, we have that for any $c > 0$ it must be $(\restrict{\mathcal{H}}{\mathcal{N}})^\rho_{\mathsf{d}(n)}(F) \geq c$ for all sufficiently large $n \gg 1$. Thus, by Lemma~\ref{lemma:Hausdorff_measure_characterization_of_S_c_n}, the pruned tree representation of $F$ itself will be an element of all the classes $\mathcal{S}^c_n$, witnessing these classes to be non-empty.   

Thus, to conclude Theorem~\ref{theorem:existence_of_compact_subset_exact_for_premeasure}, it will suffice to show the following.

\begin{proposition}
    \label{proposition:full_density_of_U_and_S}
    Let $0 < d < c,$ and let $n_0 \in \omega$ be such that $\mathcal{S}^d_{n_0}$ is non-empty. Then, for all $n \in \omega$, the set $\mathcal{U}^c_n$ is dense in $\mathcal{S}^d_{n_0}$.
\end{proposition}

This result relies on a Baire Category argument, and points to a forcing notion. 

Now, before we can prove Proposition~\ref{proposition:full_density_of_U_and_S}, let us first focus on the case when the indices of these classes agree.

\begin{lemma} 
    \label{lemma:density_of_U_and_S_same_indices}
    Let $0 < d < c$, and let $n \in \omega$ be fixed. If $\mathcal{S}^d_n$ is non-empty, then $\mathcal{U}^c_n$ is dense in $\mathcal{S}_n^d$.
\end{lemma}
\begin{proof}
    Fix a finite binary string $\zeta \in 2^{< \omega}$ having an infinite extension $Z \in \mathcal{S}^d_n$. We claim that $\zeta$ also has an extension in $\mathcal{U}^c_n \cap \mathcal{S}^d_n$. Without loss of generality, we may assume that $\zeta$ is defined for all strings of a fixed length $n' \geq n$ by extending $\zeta$ along $Z$. Call $k$ the number of length-$n'$ strings $\sigma$ for which $Z(\sigma) = 1$, and $\varepsilon := c - d > 0$. Choose $m \geq n'$ sufficiently large so that $h(\mathsf{d}(m)) < \min\seq{c / k, \varepsilon}$. And define $Y_0 \in \mathcal{P}_F$ to be an infinite binary string such that for all $\sigma \in b^{<\omega}$:
    \begin{enumerate}[label = (\roman*)]
        \item When $\operatorname{len}(\sigma) \leq n'$, set $Y_0(\sigma) = Z(\sigma)$,
        \item When $n' < \operatorname{len}(\sigma) \leq m$, set $Y_0(\sigma) = 1$ if and only if $\sigma$ is the lexicographically-least of all extensions $\nu \succeq \restrict{\sigma}{n'}$ such that $Z(\nu) = 1$, and
        \item When $\operatorname{len}(\sigma) > m$, set $Y_0(\sigma) = 1$ if and only if both $Y_0(\restrict{\sigma}{m})$ and $Z(\sigma) = 1$.
    \end{enumerate}
    
    It is straightforward to check that $Y_0$ satisfies all of the conditions for being in $\mathcal{P}_F$. Observe that by condition (ii), we will have exactly $k$ strings $\sigma$ of length $m$ for which $Y_0(\sigma) = 1$, since each of the $k$-many length-$n'$ strings represented in $Z$ will yield exactly one extension. Taking $U$ to consist of exactly these $k$ extensions of length $m$, then $\operatorname{DW}_\rho(U) = k \cdot h(\mathsf{d}(m)) < c$. This witnesses that $\phi_{n, m}(Y_0) < c$, and thus $Y_0 \in \mathcal{U}^c_n$. Note that $Y_0 \in [\zeta]$ by condition (i). 

    If $Y_0$ is also in $\mathcal{S}^d_n$, then we are done. Otherwise, let $\seq{\sigma_1, \dots , \sigma_t}$ be an enumeration of all of the strings in $b^m$ for which $Y_0(\sigma) = 0$ yet $Z(\sigma) = 1$. At least one such string exists, or else, $Y_0 = Z$, but $Z \in \mathcal{S}^d_n$. 
    
    Inductively define the infinite binary strings $Y_1, \dots, Y_t$ in $\mathcal{P}_F \cap [\zeta]$ as follows, where $0 \leq i < t$:
    \begin{itemize}
        \item If $Y_i(\sigma) = 1$, set $Y_{i + 1}(\sigma) = 1$, or
        \item If $Y_i(\sigma) = 0$, set $Y_{i + 1}(\sigma) = 1$ if and only if both $\sigma || \sigma_{i + 1}$ and $Z(\sigma) = 1$.
    \end{itemize}
    
    Again, it is straightforward to verify that these $Y_i$ are in fact elements of $\mathcal{P}_F \cap [\zeta]$. We also have that $Y_t = Z$, since any string $\sigma \in T_Z$ must be compatible either with something in $T_{Y_0}$ or with one of the strings from $\seq{\sigma_1, \dots , \sigma_t}$.
    
    We claim that $Y_i$ must be in both $\mathcal{S}_n^d$ and $\mathcal{U}^c_n$ for some $1 \leq i \leq t$. For a fixed $m' \geq m$, consider the values of $\phi_{n, m'}(Y_0), \phi_{n, m'}(Y_1), \dots, \phi_{n, m'}(Y_t)$. We have that $\phi_{n, m'}(Y_0) < c$, since $\phi_{n, m}(Y_0) < c$ and $\phi$ is monotonically-decreasing in its second coordinate. We also have that $\phi_{n, m'}(Y_t) = \phi_{n, m'}(Z) \geq d$ by assumption. Moreover, we observe for each $0 \leq i < t$ that 
    \begin{align*}
        \phi_{n, m'}(Y_i) \leq \phi_{n, m'}(Y_{i + 1}) < \phi_{n, m'}(Y_i) + \varepsilon.
    \end{align*}
    The first inequality follows from the fact that $Y_i(\sigma) = 1$ implies $Y_{i + 1}(\sigma) = 1$. So any valid cover for computing $\phi_{n, m'}(Y_{i + 1})$, after deleting any strings not present in $T_{Y_i}$, gives a cover for $Y_i$ of no greater direct $\rho$-weight, hence valid in computing $\phi_{n, m'}(Y_i)$. The second inequality follows from the fact that if $U$ is the minimal cover whose direct $\rho$-weight witnesses $\phi_{n, m'}(Y_i)$, then adding $\sigma_i$ to this cover would give a cover valid in computing $\phi_{n, m'}(Y_{i + 1})$. Thus, by the choice of $m$,
    \begin{align*}
        \phi_{n, m'}(Y_{i + 1}) \leq \phi_{n, m'}(Y_i) + h(\mathsf{d}(m)) < \phi_{n, m'}(Y_i) + \varepsilon.
    \end{align*}
    It follows that one of the values 
    \begin{align*}
        \phi_{n, m'}(Y_0), \phi_{n, m'}(Y_1), ..., \phi_{n, m'}(Y_t),
    \end{align*}
    falls in the range $[d, d + \varepsilon) = [d, c)$, as $\phi_{n, m'}(Y_0) < c$ and $\phi_{n, m'}(Y_t) \geq d$. Let $Y^{(m')}$ denote the first element $Y_i$ for which $\phi_{n, m'}(Y_i) \in [d, c)$. Since we only have finitely many $Y_i$ to consider, at least one $Y_i$ must appear as $Y^{(m')}$ for infinitely many $m' \geq m$. For such a $Y_i$, we have $Y_i \in \mathcal{U}^c_n$, as any of these $m' \geq m \geq n$ will work as a witness for $\phi_{n, m'}(Y_i) < c$. But since we can find arbitrarily large values of $m' \geq n$ for which $\phi_{n, m'}(Y_i) \geq d$, it follows by monotonicity that we in fact have $\phi_{n, m'}(Y_i) \geq d$ for all $m' \geq n$, giving $Y_i \in \mathcal{S}^d_n$, as desired. 
\end{proof}

% Really, Lemma~\ref{lemma:density_of_U_and_S_same_indices} shows that $\mathcal{U}^c_k$ is dense in $\mathcal{S}^d_n$ for all $k \leq n$ and $0 < d < c$, since $\mathcal{U}_k^c \supseteq \mathcal{U}^c_n$. 

Returning to Proposition~\ref{proposition:full_density_of_U_and_S}, by induction it now suffices to prove that if $\mathcal{U}^c_n$ is dense in $\mathcal{S}_{n_0}^d$ for some $n \geq n_0$, we also must have $\mathcal{U}_{n + 1}^c$ being dense in $\mathcal{S}_{n_0}^d$. To do this, we would want to be able to pass from a cover witnessing $Z \in \mathcal{U}^c_n$ to a cover witnessing $Z \in \mathcal{U}^c_{n + 1}$. Of course, the length-$n$ strings in that original cover will need to be replaced by proper extensions. The following definitions identify strings which can be replaced without significantly increasing the direct $\rho$-weight of the cover. 

\begin{definition}
    Let $Z \in \mathcal{P}_F$ and $\tau \in b^{<\omega}.$ Then $Z^{\tau} \in \mathcal{P}_F$ denotes the infinite binary string defined by $Z^{\tau}(\sigma) = 1$ if and only if $\sigma \mid \mid \tau$ and $Z(\sigma) = 1$. 
\end{definition}

\begin{definition}
    Let $Z \in \mathcal{P}_F$ and $n \in \omega$. Define $Z$ to be \textit{$n$-thin} if for every $\tau \in b^n$, we have
    \begin{align*}
        \lim_{m \to \infty} \phi_{n+1, m}(Z^\tau) \leq h(\mathsf{d}(n)).
    \end{align*}
    And let $\mathcal{T}_n$ denote the set of all $n$-thin sequences.
\end{definition}

As we will see, the $n$-thin sequences suffice for passing from $\mathcal{U}^c_n$ to $\mathcal{U}^c_{n+1}$. 

\begin{lemma}
    \label{lemma:use_n_thin_to_pass_n_to_n+1}
    For all $n \in \omega$ and $c > 0$, we have $\mathcal{U}^c_n \cap \mathcal{T}_n \subseteq \mathcal{U}^c_{n+1}.$
\end{lemma}
\begin{proof}
    Take $Z \in \mathcal{U}^c_n \cap \mathcal{T}_n$, and let $U = \seq{\tau_1, ..., \tau_t}$ be a minimal cover for $Z$ witnessing that $\phi_{n, m}(Z) < c$ for some $m \geq n$. 
    
    If all strings in $U$ have length at least $n + 1$, then this cover also witnesses that $\phi_{n+1, m}(Z) < c$, yielding $Z \in \mathcal{U}^c_{n + 1}$, as desired. 
    
    Otherwise, assume without loss of generality that there is some $u$ with $1 \leq u \leq t$ for which $\tau_1, \tau_2, ..., \tau_u$ are exactly the strings in $U$ of length $n$. Since $\operatorname{DW}_\rho(U) < c$, we may choose $\varepsilon > 0$ small enough so that we also have $\operatorname{DW}_\rho(U) < c - \varepsilon.$ Since $Z$ is $n$-thin, we may choose $m' \geq m$ so large that for all $\tau_i$ in $U$ of length $n$, we have
    \begin{align*}
        \phi_{n+1, m'}(Z^{\tau_i}) < h(\mathsf{d}(n)) + \frac{\varepsilon}{u}.
    \end{align*}
    For each $1 \leq i \leq t$, let $V_i$ be a minimizing cover witnessing $\phi_{n + 1, m'}(Z^{\tau_i})$, and let $U'$ be the cover obtained by replacing each string $\tau_i$ in $U$ with the strings from $V_i$. It is straightforward to verify that $U'$ is still a valid cover for computing $\phi_{n+1, m'}(Z)$. And the direct $\rho$-weight of this new cover satisfies:
    \begin{align*}
        \operatorname{DW}_\rho(U') &= \sum_{1 \leq i \leq u} \operatorname{DW}_\rho(V_i) + \operatorname{DW}_\rho(\seq{\tau_{u + 1}, ..., \tau_t})\\
        &< \sum_{1 \leq i \leq u} \p{h(\mathsf{d}(n)) + \frac{\varepsilon}{u}} + \operatorname{DW}_\rho(\seq{\tau_{u + 1}, ..., \tau_t}) = \operatorname{DW}_\rho(U) + \varepsilon < c.
    \end{align*}
    This proves that $Z$ is an element of $\mathcal{U}_{n+1}^c$, as desired.
\end{proof}

It remains to show that the $n$-thin elements are sufficiently numerous within each non-empty set $\mathcal{S}^d_{n_0}.$
\begin{lemma}
    \label{lemma:n_thins_dense_in_complement}
    Let $d > 0$ and let $n_0 \in \omega$ be such that $\mathcal{S}^d_{n_0}$ is non-empty. Then for all $n \geq n_0$, the set $\mathcal{T}_n$ is dense in $\mathcal{S}^d_{n_0}$.
\end{lemma}
\begin{proof}
    Fix $n \geq n_0$ and a finite binary string $\zeta \in 2^{< \omega}$ having an infinite extension $Z \in \mathcal{S}^d_{n_0}$. We claim that $\zeta$ has an extension in $\mathcal{T}_n \cap \mathcal{S}^d_{n_0}$. Without loss of generality, we may assume that $\zeta$ is defined for all strings of length $n$ by extending $\zeta$ along $Z$. For each $\tau \in b^n$, denote by $\zeta^\tau \in 2^{\operatorname{len}(\zeta)}$ the string defined by
    \begin{align*}
        \zeta^\tau(\sigma) = 1 \iff \sigma \mid \mid \tau \text{ and } \zeta(\sigma) = 1.
    \end{align*}

    We start by defining an element $Y \in \mathcal{T}_n \cap [\zeta]$, and later check that $Y \in \mathcal{S}^d_{n_0}$, witnessing $\mathcal{T}_n$ being dense in $\mathcal{S}^d_{n_0}$.
    
    To define the element $Y \in \mathcal{P}_F$, it suffices to both specify $Y^\tau$ for each $\tau \in b^n$ and further ensure that each $Y^\tau$ extends $\zeta^\tau$. In this way, the final sequence $Y$ will be well-defined.
    
    If a string $\sigma \in b^{<\omega}$ is compatible with two distinct strings $\tau_1, \tau_2 \in b^n$, then $\sigma$ must be a common prefix for both of these strings, implying $\operatorname{len}(\sigma) < n$. Therefore, $Y(\sigma)$ will be determined by $\zeta(\sigma)$ by the assumption that $\zeta$ is defined on all strings of length up to $n$. It is also straightforward to show that $Y$ will be an element of $\mathcal{P}_F$ so long as each $Y^\tau$ is.
    
    Given $\tau \in b^n$, we say that $\tau$ is \textit{$n$-thin for $Z$} if:    
    \begin{align*}
        \lim_{m \to \infty} \phi_{n + 1, m}(Z^\tau) \leq h(\mathsf{d}(n)).
    \end{align*}
    We further collect into $\mathcal{T}^Z_n \subseteq b^n$ all the $n$-thin strings for $Z$. Now, define $Y^\tau = Z^\tau$ whenever $\tau \in \mathcal{T}^Z_n$. Noting that $\zeta^\tau \preceq Z^\tau$, the claim is satisfied for these $\tau$. 
    
    Otherwise, $\tau \not\in \mathcal{T}^Z_n$, so it must be that
    \begin{align*}
        \lim_{m \to \infty} \phi_{n + 1, m}(Z^\tau) > h(\mathsf{d}(n)),
    \end{align*}
    which, by monotonicity, implies that $\phi_{n + 1, m}(Z^\tau) > h(\mathsf{d}(n))$ for all $m \geq n + 1$, hence $Z^\tau \in \mathcal{S}_{n + 1}^{h(\mathsf{d}(n))}.$
    
    Take $(r_i)_{i \in \omega}$ to be any decreasing sequence of positive real numbers converging to $0$. By Lemma~\ref{lemma:density_of_U_and_S_same_indices}, it follows that $\mathcal{U}_{n + 1}^{h(\mathsf{d}(n)) + r_i}$ is dense in $\mathcal{S}_{n + 1}^{h(\mathsf{d}(n))}$ for each $i \in \omega$. And by the Baire Category Theorem, the intersection $\bigcap_i \mathcal{U}_{n + 1}^{h(\mathsf{d}(n)) + r_i}$ is also dense in $\mathcal{S}_{n + 1}^{h(\mathsf{d}(n))}$. And since $\zeta^\tau \preceq Z^\tau \in \mathcal{S}_{n + 1}^{h(\mathsf{d}(n))}$, it follows that $[\zeta^\tau] \cap \mathcal{S}_{n + 1}^{h(\mathsf{d}(n))} \neq \emptyset$. Thus, we may fix a sequence $Y^\tau \in [\zeta^\tau] \cap  \mathcal{S}_{n + 1}^{h(\mathsf{d}(n))} \cap \bigcap_i \mathcal{U}_{n + 1}^{h(\mathsf{d}(n)) + r_i}$. 

    The fact that $Y^\tau \in \mathcal{S}_{n + 1}^{h(\mathsf{d}(n))}$ implies that $\phi_{n + 1, m}(Y^\tau) \geq h(\mathsf{d}(n))$ for all $m \geq n + 1$. Yet, for any $\varepsilon > 0$, one may choose $i \gg 1$ for which $r_i < \varepsilon$. Then, since $Y^\tau \in \mathcal{U}_{n + 1}^{h(\mathsf{d}(n)) + r_i}$, there exists some $m \geq n + 1$ for which 
    \begin{align*}
        \phi_{n + 1, m}(Y^\tau) < h(\mathsf{d}(n)) + r_i < h(\mathsf{d}(n)) + \varepsilon.
    \end{align*}
    So, by the monotonicity of $\phi_{n + 1, m}(Y^\tau)$ in its second component, we conclude that
    \begin{align*}
        \lim_{m \to \infty} \phi_{n + 1, m}(Y^\tau) = h(\mathsf{d}(n)).
    \end{align*}
    This construction guarantees that $Y$ both extends $\zeta$ and is $n$-thin.

    Now we show that $Y$ is indeed an element of $\mathcal{S}_{n_0}^d$, i.e., that $\phi_{n_0,  m}(Y) \geq d$ for each $m \geq n_0$. By monotonicity, it suffices to only consider $m \geq n \geq n_0$. Fix such an $m \geq n$ and suppose $U \subseteq b^{<\omega}$ is a cover valid for computing $\phi_{n_0, m}(Y)$. We wish to show that $\operatorname{DW}_\rho(U) \geq d$. 

    Without loss of generality, we may take $U$ to be prefix-free, for any strings in $U$ having proper prefixes in $U$ may be deleted while maintaining the cover's validity in computing $\phi_{n_0, m}(Y)$ and not increasing the direct $\rho$-weight of the resulting cover. We plan to push the cover $U$ to a cover $U'$ valid instead for computing $\phi_{n_0, m}(Z)$ and satisfying $\operatorname{DW}_\rho(U') \leq \operatorname{DW}_\rho(U)$. Since $Z$ is an element of $\mathcal{S}_{n_0}^d$, this would yield $\operatorname{DW}_\rho(U) \geq \operatorname{DW}_\rho(U') \geq d$, putting $Y \in \mathcal{S}_{n_0}^d$, as $U$ was arbitrary.
    
    Consider all $\tau \in b^n \setminus \mathcal{T}^Z_n$. Denote by $V_\tau$ the collection of all proper extensions of $\tau$ in $U$. If $V_\tau = \emptyset$, do nothing. Otherwise, replace in the cover $U$ all strings in $V_\tau$ by $\tau$ itself. Let us argue why this replacement will not increase the direct $\rho$-weight of the cover. 
    
    Note that all strings in $V_\tau$ have length between $n + 1$ and $m$, inclusive. And if $\sigma \in b^m$ is such that $Y^\tau(\sigma) = 1$, then $\sigma$ must extend some element of $V_\tau$. Given that the original set $U$ is a cover for $Y$, we have that $\sigma$ extends some element $\mu \in U$ compatible with $\tau$. If $\mu \preceq \tau$, then $\mu$ would be a proper prefix of any element in $V_\tau$, contradicting the fact that $V_\tau$ was non-empty and $U$ was prefix-free. Therefore, it must be that $\mu \succ \tau$, meaning $\mu \in V_\tau.$ Taken together, we see that $V_\tau$ is valid cover for computing $\phi_{n + 1, m}(Y^\tau)$. Since $\tau \not\in \mathcal{T}^Z_n$, the construction of $Y$ yields that $\phi_{n + 1, m}(Y^\tau) \geq h(\mathsf{d}(n))$, so $\operatorname{DW}_\rho(V_\tau) \geq h(\mathsf{d}(n)).$ So, as $\tau$ has length $n$, replacing all of $V_\tau$ with $\tau$ cannot increase the direct $\rho$-weight of $U$.
    
    Let $U'$ be the cover obtained by performing this replacement for all $\tau \in b^n \setminus \mathcal{T}^Z_n$. It remains to check that $U'$ valid for computing $\phi_{n_0, m}(Z)$. Since the replacement process may only add strings of length $n$, and $n_0 \leq n \leq m$, the sequence $U'$ satisfies the length condition. 
    
    For the covering condition, suppose we have a string $\sigma \in b^m$ for which $Z(\sigma) = 1$. Let $\tau = \restrict{\sigma}{n}$. Suppose $\tau \in \mathcal{T}^Z_n$. Then by construction, we have $Y^\tau = Z^\tau$, so $Y^\tau(\sigma) = Z^\tau(\sigma) = 1$. As the original cover $U$ is valid for $\phi_{n_0, m}(Y)$, this implies that $U$ contains some prefix of $\sigma$ compatible with $\tau$. This prefix will still be in $U'$, since we only removed strings which extended elements of $b^n \setminus \mathcal{T}^Z_n$.
    
    Otherwise, $\tau \notin \mathcal{T}^Z_n$. The original cover $U$ is valid for computing $\phi_{n_0, m}(Y)$, and hence valid for $\phi_{n_0, m}(Y^\tau)$, as well. Given that $\phi_{n + 1, m}(Y^\tau) \geq h(\mathsf{d}(n))$, there must be at least one string $\sigma' \in b^m$ such that $Y^\tau(\sigma') = 1$. So, this $\sigma'$ has a prefix $\mu \in U$ compatible with $\tau$. If $\mu \preceq \sigma$, then $\mu$ is also a prefix of $\sigma$ and is not removed when passing to the cover $U'$. Otherwise, $\mu$ \textit{is} removed and replaced with $\tau$. Either way, we get that $U'$ contains a prefix of $\sigma$, completing the argument that $U'$ is a valid cover for computing $\phi_{n_0, m}(Z)$, as desired. 
\end{proof}

Now we may prove the key density result.

\begin{proof}[Proof of Proposition~\ref{proposition:full_density_of_U_and_S}]
    By Lemma~\ref{lemma:density_of_U_and_S_same_indices} and since the sequence $(\mathcal{U}^c_n)_n$ is nested-decreasing, we get that $\mathcal{U}^c_n$ is dense in $\mathcal{S}^d_{n_0}$ for all $n \leq n_0$. 

    Assume we have shown that $\mathcal{U}^c_n$ is dense in $\mathcal{S}^d_{n_0}$ for some $n \geq n_0$. We now show that $\mathcal{U}^c_{n + 1}$ is dense in $\mathcal{S}^d_{n_0}$ as well. 

    Fix a finite binary string $\zeta \in 2^{<\omega}$ having an infinite extension $Z \in \mathcal{S}^d_{n_0} \cap \mathcal{U}^c_n$. Since the set $\mathcal{U}^c_n$ is open, we may choose an extension $\zeta \preceq \zeta' \prec Z$ for which $[\zeta'] \cap \mathcal{S}^d_{n_0} \subseteq \mathcal{U}^c_n$. Now, as $[\zeta'] \cap \mathcal{S}^d_{n_0} \neq \emptyset$ ($Z$ is contained in their intersection), Lemma~\ref{lemma:n_thins_dense_in_complement} implies there exists an $n$-thin sequence $Y \in [\zeta'] \cap \mathcal{S}^d_{n_0} \cap \mathcal{T}_n$. And by the choice of $\zeta'$, we have $Y \in \mathcal{T}_n \cap \mathcal{U}^c_n$, so by Lemma~\ref{lemma:use_n_thin_to_pass_n_to_n+1}, $Y \in \mathcal{U}^c_{n + 1}$. This shows that $\mathcal{U}^c_{n + 1}$ is dense in $\mathcal{S}^d_{n_0}$. The result now follows by induction.    
\end{proof}

% \section{Miscellaneous}

% \subsection{Correspondence Principle for Measure PTS}

% {\color{red} If going to include, need more details showing double-jump suffices.}
% % Note that certain Correspondence Principles are true for them too!

% We note that it is possible to verify in the proof of the formula for $\log \mathcal{H}^s(X)$ that it admits a correspondence principle of its own for effectively closed sets. If $X$ is an effectively closed set relative to some oracle $B$, and $s \geq 0$ is left-c.e. in $B$, we may use $B''$ (the double Turing-jump of $B$) to compute $\mathcal{H}^s(X)$ as a real number. This is needed in the construction of a strong Solovay-$s$-test covering $X$ of appropriate prefix $s$-weight.

% \begin{proposition}
%     \label{proposition:correspondence_principle_log_s_measure_and_null_and_finite}
%     For every $B \in 2^{\leq \omega}$, $\Pi^{0, B}_1$-class $X \subseteq 2^\omega$ and $B$-left-c.e. $s \geq 0$,
%     \begin{align*}
%         \log \mathcal{H}^s(X) &=^+ \sup_{x \in X} \liminf_{r \to \infty}  \operatorname{\Delta M}^{B''}_s(\restrict{x}{r}).
%     \end{align*}
%     In particular,
%     \begin{align*}
%         \mathcal{H}^s(X) > 0 & \iff (\exists x \in X) \liminf_{r \to \infty} \operatorname{\Delta M}^{B''}_s(\restrict{x}{r}) > -\infty, \quad \text{and}\\         
%         \mathcal{H}^s(X) < \infty &\iff \sup_{x \in X} \liminf_{r \to \infty} \operatorname{\Delta M}^{B''}_s(\restrict{x}{r}) < \infty.
%     \end{align*} 
% \end{proposition}
% It is not clear that the characterization of $\sigma$-finiteness with respect to $s$-measure admits a correspondence principle for classes in some level of the lightface Borel hierarchy. This is because $\sigma$-finiteness for a (definably) simple set does not necessarily mean the set is covered by a countable collection of simple, finite-measure sets.